本实验的攻击算法按照实验指导书中的思路设计，并与 \verb|attack\_slow.py| 对应。第 3 节已经说明了 Padding Oracle 攻击的原理，本节整理算法的执行流程。

\subsection{整体攻击思路}

程序首先将目标密文拆分为
\[
    IV,\ C_1,\ C_2,\ C_3
\]
四个分组，然后从最后一个目标分组开始，逐块恢复对应的中间状态和明文。

\subsection{算法流程}

\begin{enumerate}
    \item 取该目标分组真实前一密文分组作为初始参考分组（如果目标分组不是最后一块，需要穷搜前一块的最后一个字节）；
    \item 通过逐个修改参考分组的字节，判断原始明文末尾的填充长度；
    \item 根据已知填充长度，先恢复中间状态末尾字节；
    \item 修改前一块的字节，构造新的填充值；
    \item 穷搜恢复前一个中间状态字节；
    \item 当整个中间状态恢复完成后，再与真实前一密文分组异或，得到目标明文分组。
\end{enumerate}

这一设计与实验指导书中的推导是一致的，也和第 3 节中的原理说明完全对应。

\subsection{按块恢复的执行过程}

对每个目标分组，程序都会先获得该分组对应的中间状态后缀，再逐步向前扩展。这个过程与原理部分的例子一致：先根据服务器反馈判断原始填充长度，再利用已知后缀去构造新的合法填充，逐字节恢复更多中间状态。等到整个中间状态都恢复完成后，再与真实前一密文分组异或，即可得到当前明文分组。

\subsection{实际提交密文格式}

理论上，Padding Oracle 攻击只需要最后两个分组。但在本实验环境中，服务器要求提交完整密文；实际实现时，只需把前面的无关分组设为 0 即可，对攻击本身没有影响。

\subsection{算法复杂度}

对每个待恢复字节而言，最坏情况下需要枚举 256 种可能取值。每个 AES 分组有 16 个字节，因此恢复一个分组最多需要
\[
    16 \times 256 = 4096
\]
次 oracle 查询。
本实验目标密文除 $IV$ 外共有 3 个数据分组，因此理论上的查询上界为
\[
    3 \times 16 \times 256 = 12288.
\]
实际运行时，并不是每个字节都会枚举满 256 次，因此真实查询次数通常小于该上界。与此同时，程序是按块逐步恢复的，因此一旦某个分组恢复完成，就可以立即得到对应的明文结果，便于随时检查程序是否运行正常。
